% Scope: Establish exact frame transformations and RF phase conventions before treating excitation, inversion and selective pulses.
% Status: architecture only; no chapter prose or completed problems yet.
% Prerequisites: Chapters 2 and 3.
% Planned worked examples: Rotating-frame transformation and flip angle of a rectangular RF pulse.
% Planned figures: Laboratory and rotating axes; effective-field trajectories.
% Planned challenge problems: Derive a detuned pulse propagator; debug a factor-of-two RF amplitude error.
\chapter{RF Excitation and the Rotating Frame}
\label{ch:04}

\section{Rotating coordinate systems}
\label{sec:04:rotating-coordinate-systems}
\subsection{Active rotations versus passive coordinates}
\subsection{Transformation of time derivatives}
\subsection{Detuning and the effective field}

\section{Circular RF components}
\label{sec:04:circular-rf-components}
\subsection{Real fields and complex phasors}
\subsection{Co-rotating B1-plus and counter-rotating fields}
\subsection{Linear polarization and the rotating-wave approximation}

\section{Hard-pulse dynamics}
\label{sec:04:hard-pulse-dynamics}
\subsection{Flip angle, RF phase and rotation axis}
\subsection{On-resonance excitation and inversion}
\subsection{Off-resonance excitation and bandwidth}

\section{Pulse imperfections and calibration}
\label{sec:04:pulse-imperfections-and-calibration}
\subsection{Transmit inhomogeneity and flip-angle errors}
\subsection{Finite-duration effects and relaxation}
\subsection{Phase cycling and receiver phase}

\section{Selective excitation as an introduction}
\label{sec:04:selective-excitation-as-an-introduction}
\subsection{Slice-selection gradients and RF bandwidth}
\subsection{Small-tip response and pulse envelopes}
\subsection{Boundaries of the hard-pulse approximation}

\section{Worked examples}
\label{sec:04:worked-examples}
% Planned calculations: Rotating-frame transformation and flip angle of a rectangular RF pulse.
\subsection{Derivation and quantitative calculation}
\subsection{Assumptions, limiting cases and scanner interpretation}

\section{Challenge problems}
\label{sec:04:challenge-problems}
% Problem design: Derive a detuned pulse propagator; debug a factor-of-two RF amplitude error.
% Use challengeproblem and stable labels prob:04:<topic>.
% Complete solutions belong in Appendix E, section for this chapter.
% Use \begin{solution}{prob:04:<topic>} ... \end{solution} there.
\subsection{Derivation and quantitative design}
\subsection{Model criticism and integrated reasoning}
